\documentclass[../Project.tex]{subfiles}
\begin{document}

Chương này trình bày chi tiết về quá trình phân tích thiết kế, xây dựng ứng dụng, kiểm thử và triển khai thực tế hệ thống trò chơi trực tuyến nhiều người chơi "The First Sect Origin". Nội dung bao gồm việc lựa chọn và đặc tả thiết kế kiến trúc phần mềm, mô tả cấu trúc các gói, thiết kế các lớp đối tượng chính, vẽ sơ đồ thực thể liên kết cơ sở dữ liệu, liệt kê chi tiết các công cụ và thư viện phát triển, trình bày các kịch bản kiểm thử chức năng và cuối cùng là phương án triển khai hệ thống máy chủ.

\section{Thiết kế kiến trúc}
\label{sec:thiet_ke_kien_truc}

\subsection{Lựa chọn kiến trúc phần mềm}
\label{sec:lua_chon_kien_truc}

Trong việc phát triển hệ thống trò chơi trực tuyến nhiều người chơi đòi hỏi khả năng phản hồi nhanh và chịu tải lớn, thiết kế kiến trúc đóng vai trò quyết định độ ổn định lâu dài của sản phẩm. Hệ thống "The First Sect Origin" được thiết kế dựa trên sự kết hợp giữa kiến trúc \textbf{Vi dịch vụ (Microservices)} ở mức độ hệ thống tổng quát và kiến trúc \textbf{Clean Architecture} ở mức độ cấu trúc mã nguồn bên trong từng dịch vụ.

\subsubsection{Mô hình kiến trúc Vi dịch vụ (Microservices)}
Hệ thống máy chủ backend được chia tách thành các dịch vụ độc lập, chuyên biệt hóa nhiệm vụ và giao tiếp với nhau thông qua giao thức gRPC hiệu năng cao:
\begin{itemize}
    \item \textbf{Gateway Service}: Đóng vai trò là điểm tiếp nhận yêu cầu (Single Entry Point) duy nhất từ phía client. Dịch vụ này thực hiện chức năng xác thực mã thông báo JWT của người dùng, giới hạn tần suất yêu cầu (Rate Limiting) nhằm chống tấn công từ chối dịch vụ, và điều phối các cuộc gọi gRPC tới các vi dịch vụ tương ứng ở mạng nội bộ. Đồng thời, Gateway còn chứa logic "Roll Server" để điều hướng người chơi kết nối vào phân vùng (Zone) thích hợp khi hệ thống mở rộng.
    \item \textbf{Login Server}: Chịu trách nhiệm xử lý toàn bộ các yêu cầu liên quan đến tài khoản người dùng, bao gồm đăng ký, đăng nhập, liên kết tài khoản mạng xã hội (Google, Apple, Facebook) và cấp phát token xác thực. Dịch vụ này truy cập trực tiếp vào cơ sở dữ liệu toàn cục (Global Database) chứa thông tin tài khoản dùng chung giữa các Zone.
    \item \textbf{World Server}: Quản lý toàn bộ trạng thái thế giới trò chơi và dữ liệu phiêu lưu của người chơi thuộc một phân vùng cụ thể. Các logic cốt lõi như xây dựng và nâng cấp tông môn, sản sinh tài nguyên (Gỗ, Đá, Vàng (Item ID: 00001)), quản lý đệ tử, sắp xếp đội hình chiến đấu, thực hiện gacha chiêu mộ và theo dõi độ kiên nhẫn (pity) đều diễn ra tại đây. World Server lưu trữ trạng thái bền vững trong Game Database của riêng phân vùng đó.
    \item \textbf{Combat Server}: Dịch vụ tính toán phi trạng thái (Stateless) chuyên xử lý và mô phỏng các trận đấu thẻ tướng chiến thuật 2D. Khi người chơi tham gia vượt ải, World Server sẽ gửi thông tin đội hình người chơi và đội hình quái vật sang Combat Server qua gRPC. Combat Server thực hiện mô phỏng toàn bộ tiến trình trận đấu theo lượt, tính toán sát thương, hiệu ứng và trả về kết quả thắng/thua cùng nhật ký trận đấu chi tiết (battle log).
\end{itemize}

\subsubsection{Mô hình Clean Architecture cho từng dịch vụ}
Mỗi dịch vụ được tổ chức mã nguồn theo 4 lớp cốt lõi của Clean Architecture nhằm giảm thiểu sự phụ thuộc trực tiếp vào framework hoặc công cụ lưu trữ dữ liệu, tăng khả năng viết kiểm thử đơn vị độc lập:
\begin{enumerate}
    \item \textbf{Domain Entities Layer}: Chứa các định nghĩa cấu trúc dữ liệu thuần túy (như \texttt{Player}, \texttt{PlayerHero}, \texttt{PlayerBuilding}) đại diện cho thực thể nghiệp vụ. Lớp này không phụ thuộc vào bất kỳ thư viện ngoài hay hệ quản trị cơ sở dữ liệu nào.
    \item \textbf{Use Cases / Services Layer}: Chứa toàn bộ quy tắc nghiệp vụ (Business Rules) của game. Ví dụ, dịch vụ \texttt{BuildingService} chứa logic kiểm tra điều kiện nâng cấp công trình (người chơi có đủ tài nguyên Gỗ, Đá, Vàng (Item ID: 00001) không, cấp độ của Điện chính là bao nhiêu) và thực hiện trừ tài nguyên trước khi bắt đầu nâng cấp.
    \item \textbf{Interface Adapters (Handlers/Controllers) Layer}: Đóng vai trò cầu nối chuyển đổi dữ liệu. Trong dự án này, lớp này là các gRPC handlers nhận các đối tượng request dạng Protobuf (ví dụ \texttt{UpgradeBuildingRequest}), chuyển đổi sang dạng struct Go, gọi Use Case tương ứng và sau đó chuyển đổi kết quả thành Protobuf response trả về cho client.
    \item \textbf{Infrastructure (Repositories/Drivers) Layer}: Chứa mã nguồn thao tác với cơ sở dữ liệu (PostgreSQL qua thư viện \texttt{pgx}, Redis qua \texttt{go-redis}) và giao tiếp ngoại vi. Lớp này hiện thực hóa các giao diện (Interface) được định nghĩa ở lớp nghiệp vụ thông qua cơ chế tiêm phụ thuộc (Dependency Injection).
\end{enumerate}

\subsection{Thiết kế tổng quan}
\label{sec:thiet_ke_tong_quan}

Dưới đây trình bày sơ đồ gói UML (UML Package Diagram) thể hiện mối quan hệ phụ thuộc giữa các thành phần phần mềm trong toàn bộ hệ thống. Các gói được phân chia thành các tầng rõ ràng theo quy tắc thiết kế kiến trúc: tầng trên phụ thuộc tầng dưới, các gói ở cùng một tầng nghiệp vụ không có sự phụ thuộc chéo lẫn nhau, và các giao tiếp mạng được chuẩn hóa qua gRPC.

\subsection{Hiện thực Giao diện Người dùng}
\label{sec:hien_thuc_giao_dien}

\subsubsection{Giao diện Quản lý Tông môn}

\textbf{1. Màn hình Xây dựng Công trình:}

Giao diện trực quan trên lưới (grid) cho phép người chơi đặt nền móng và tiến hành xây dựng các tòa nhà mới trong Tông môn.

\begin{figure}[H]
    \centering
    \includegraphics[width=0.85\textwidth]{Image/ui_building_construct.png}
    \caption{Màn hình Xây dựng Công trình}
    \label{fig:ui_building_construct}
\end{figure}

\textbf{2. Màn hình Menu Xây dựng:}

Cung cấp danh sách các công trình kiến trúc có thể xây dựng. Hiển thị thông tin tổng quan, thời gian hoàn thành và chi phí tài nguyên yêu cầu cho từng loại công trình.

\begin{figure}[H]
    \centering
    \includegraphics[width=0.85\textwidth]{Image/ui_build_menu.png}
    \caption{Màn hình Menu Xây dựng}
    \label{fig:ui_build_menu}
\end{figure}

\textbf{3. Màn hình Thông tin Công trình:}

Hiển thị thông tin chi tiết của một công trình đã xây dựng, bao gồm cấp độ hiện tại, tiến độ nâng cấp, và các chỉ số tài nguyên/sức mạnh cộng thêm cho Tông môn.

\begin{figure}[H]
    \centering
    \includegraphics[width=0.85\textwidth]{Image/ui_building_info.png}
    \caption{Màn hình Thông tin Công trình}
    \label{fig:ui_building_info}
\end{figure}

\subsubsection{Giao diện Đệ tử và Túi đồ}

\textbf{1. Màn hình Danh sách Đệ tử:}

Giao diện hiển thị và quản lý toàn bộ Đệ tử người chơi đang sở hữu. Người chơi có thể phân loại, xem thông tin tóm tắt và lựa chọn đệ tử để bồi dưỡng.

\begin{figure}[H]
    \centering
    \includegraphics[width=0.85\textwidth]{Image/ui_disciple_list.png}
    \caption{Màn hình Danh sách Đệ tử}
    \label{fig:ui_disciple_list}
\end{figure}

\textbf{2. Màn hình Sử dụng Vật phẩm (Item):}

Giao diện túi đồ (Inventory) nơi người chơi kiểm tra tài nguyên, xem chi tiết và sử dụng các loại vật phẩm (đan dược, rương tài nguyên) để hỗ trợ quá trình tu luyện.

\begin{figure}[H]
    \centering
    \includegraphics[width=0.85\textwidth]{Image/ui_use_item.png}
    \caption{Màn hình Sử dụng Vật phẩm}
    \label{fig:ui_use_item}
\end{figure}

\subsubsection{Giao diện Bản đồ và Kết quả Chiến đấu}

\textbf{1. Màn hình Chọn Màn chơi:}

Bản đồ chiến dịch nơi người chơi lựa chọn các ải (Level) để tiến hành khiêu chiến. Giao diện hiển thị tuyến đường, số sao đạt được và phần thưởng rớt ra ở từng ải.

\begin{figure}[H]
    \centering
    \includegraphics[width=0.85\textwidth]{Image/ui_level_select.png}
    \caption{Màn hình Chọn Màn chơi}
    \label{fig:ui_level_select}
\end{figure}

\textbf{2. Màn hình Kết quả - Thắng (Victory):}

Giao diện tổng kết sau khi người chơi vượt ải thành công, hiển thị lượng kinh nghiệm (EXP), vàng và các tài nguyên (gỗ, đá, vàng (Item ID: 00001)) thu thập được từ trận đấu.

\begin{figure}[H]
    \centering
    \includegraphics[width=0.85\textwidth]{Image/ui_victory.png}
    \caption{Màn hình Kết quả Trận đấu (Thắng)}
    \label{fig:ui_victory}
\end{figure}

\textbf{3. Màn hình Kết quả - Thua (Defeat):}

Giao diện thông báo kết quả khi đội hình người chơi bị hạ gục. Giao diện cung cấp tùy chọn thoát hoặc thử lại, đồng thời gợi ý nâng cấp sức mạnh.

\begin{figure}[H]
    \centering
    \includegraphics[width=0.85\textwidth]{Image/ui_defeat.png}
    \caption{Màn hình Kết quả Trận đấu (Thua)}
    \label{fig:ui_defeat}
\end{figure}

\section{Kiểm thử}
\label{sec:kiem_thu}

Quá trình kiểm thử phần mềm áp dụng hai kỹ thuật chính: \textbf{Phân hoạch tương đương (Equivalence Partitioning)} và \textbf{Phân tích giá trị biên (Boundary Value Analysis)} để thiết kế các trường hợp kiểm thử cho các chức năng quan trọng nhất, đảm bảo tính đúng đắn của nghiệp vụ game và phòng tránh các lỗ hổng gian lận tài sản.

\subsubsection{Các kịch bản kiểm thử chi tiết}

\paragraph{Kịch bản 1: Kiểm thử chức năng Đăng nhập hệ thống}

\begin{table}[H]
\centering
\begin{tabular}{|l|p{10.5cm}|}
\hline
\textbf{Thuộc tính} & \textbf{Nội dung mô tả} \\ \hline
\textbf{Mã kịch bản} & TC-AUTH-01 \\ \hline
\textbf{Tên kịch bản} & Đăng nhập thành công vào hệ thống \\ \hline
\textbf{Mục tiêu} & Xác minh người chơi có thể đăng nhập thành công với tài khoản và mật khẩu hợp lệ. \\ \hline
\textbf{Điều kiện} & Tài khoản đã được đăng ký trước trên hệ thống và đang hoạt động. \\ \hline
\textbf{Các bước thực hiện} & 
1. Mở ứng dụng client, truy cập giao diện Đăng nhập. \\
& 2. Nhập tên tài khoản hợp lệ vào ô tài khoản. \\
& 3. Nhập mật khẩu chính xác vào ô mật khẩu. \\
& 4. Nhấn nút "Đăng nhập". \\ \hline
\textbf{Dữ liệu đầu vào} & Username: \texttt{123456}, Password: \texttt{123456} \\ \hline
\textbf{Kết quả mong đợi} & Hệ thống xác thực thành công, sinh mã JWT token hợp lệ gửi về client, giao diện chuyển sang màn hình chính. \\ \hline
\textbf{Kết quả thực tế} & Đăng nhập thành công, nhận JWT token, client tự động lưu token và chuyển cảnh vào game. \\ \hline
\textbf{Trạng thái} & \textbf{PASS} \\ \hline
\end{tabular}
\caption{Kịch bản kiểm thử TC-AUTH-01: Đăng nhập thành công}
\label{tab:tc_auth_01}
\end{table}

\begin{table}[H]
\centering
\begin{tabular}{|l|p{10.5cm}|}
\hline
\textbf{Thuộc tính} & \textbf{Nội dung mô tả} \\ \hline
\textbf{Mã kịch bản} & TC-AUTH-02 \\ \hline
\textbf{Tên kịch bản} & Đăng nhập thất bại do mật khẩu không chính xác \\ \hline
\textbf{Mục tiêu} & Đảm bảo hệ thống từ chối các yêu cầu đăng nhập sử dụng mật khẩu sai. \\ \hline
\textbf{Điều kiện} & Tài khoản đã được đăng ký trước trên hệ thống. \\ \hline
\textbf{Các bước thực hiện} & 
1. Mở ứng dụng client, truy cập giao diện Đăng nhập. \\
& 2. Nhập tên tài khoản hợp lệ vào ô tài khoản. \\
& 3. Nhập mật khẩu không chính xác vào ô mật khẩu. \\
& 4. Nhấn nút "Đăng nhập". \\ \hline
\textbf{Dữ liệu đầu vào} & Username: \texttt{123456}, Password: \texttt{WrongPass1} \\ \hline
\textbf{Kết quả mong đợi} & Hệ thống từ chối xác thực, trả về mã lỗi bảo mật, client hiển thị thông báo "Mật khẩu không chính xác". \\ \hline
\textbf{Kết quả thực tế} & Đăng nhập thất bại, nhận thông báo lỗi từ server, client giữ nguyên giao diện đăng nhập và thông báo sai mật khẩu. \\ \hline
\textbf{Trạng thái} & \textbf{PASS} \\ \hline
\end{tabular}
\caption{Kịch bản kiểm thử TC-AUTH-02: Đăng nhập thất bại do sai mật khẩu}
\label{tab:tc_auth_02}
\end{table}

\paragraph{Kịch bản 2: Kiểm thử chức năng Thu thập tài nguyên công trình}
Kịch bản kiểm thử đảm bảo việc cộng tài sản chính xác và không cho phép người chơi khai thác vượt sản lượng sản sinh thực tế:

\begin{table}[H]
\centering
\begin{tabular}{|l|p{10.5cm}|}
\hline
\textbf{Thuộc tính} & \textbf{Nội dung mô tả} \\ \hline
\textbf{Mã kịch bản} & TC-COLL-01 \\ \hline
\textbf{Tên kịch bản} & Thu hoạch tài nguyên thành công từ công trình sản xuất \\ \hline
\textbf{Mục tiêu} & Xác minh tài sản của người chơi được cộng chính xác dựa trên sản lượng tích lũy thực tế của công trình. \\ \hline
\textbf{Điều kiện} & Người chơi đang ở màn hình Tông môn, công trình (Linh điền) thuộc sở hữu của người chơi và thời gian tích lũy sản lượng lớn hơn 0. \\ \hline
\textbf{Các bước thực hiện} & 
1. Người chơi bấm chọn biểu tượng vàng (Item ID: 00001) nổi lên phía trên Linh điền. \\
& 2. Client gửi yêu cầu thu hoạch \texttt{CollectResourcesRequest} chứa ID công trình lên World Server qua gRPC. \\ \hline
\textbf{Dữ liệu đầu vào} & ID công trình: \texttt{1024}, Thời gian tích lũy: 3600 giây (hiệu suất 100 vàng (Item ID: 00001)/giờ). \\ \hline
\textbf{Kết quả mong đợi} & World Server xác nhận thông tin, tính toán cộng thêm 100 vàng (Item ID: 00001) vào tài khoản của người chơi, cập nhật cơ sở dữ liệu và trả về phản hồi thành công. \\ \hline
\textbf{Kết quả thực tế} & Thu hoạch thành công, số dư vàng (Item ID: 00001) của người chơi trên thanh trạng thái tăng thêm 100 đơn vị, cơ sở dữ liệu ghi nhận chính xác. \\ \hline
\textbf{Trạng thái} & \textbf{PASS} \\ \hline
\end{tabular}
\caption{Kịch bản kiểm thử TC-COLL-01: Thu hoạch tài nguyên thành công}
\label{tab:tc_collect_01}
\end{table}

\begin{table}[H]
\centering
\begin{tabular}{|l|p{10.5cm}|}
\hline
\textbf{Thuộc tính} & \textbf{Nội dung mô tả} \\ \hline
\textbf{Mã kịch bản} & TC-COLL-02 \\ \hline
\textbf{Tên kịch bản} & Ngăn chặn thu hoạch tài nguyên khi bấm liên tục (Spam) \\ \hline
\textbf{Mục tiêu} & Ngăn chặn việc gửi liên tiếp các yêu cầu thu hoạch nhằm mục đích gian lận hoặc làm quá tải máy chủ. \\ \hline
\textbf{Điều kiện} & Người chơi vừa thực hiện thu hoạch tài nguyên thành công. \\ \hline
\textbf{Các bước thực hiện} & 
1. Người chơi thực hiện click liên tiếp 2 lần trong vòng dưới 2 giây vào biểu tượng thu hoạch trên công trình. \\
& 2. Client gửi liên tục 2 yêu cầu gRPC thu hoạch lên World Server. \\ \hline
\textbf{Dữ liệu đầu vào} & ID công trình: \texttt{1024}, Khoảng cách giữa 2 yêu cầu: 0,5 giây. \\ \hline
\textbf{Kết quả mong đợi} & Yêu cầu thứ nhất xử lý thành công. Yêu cầu thứ hai bị máy chủ từ chối, trả về mã lỗi hoặc thông báo sản lượng thu hoạch bằng 0 do chưa đủ thời gian chờ tối thiểu. \\ \hline
\textbf{Kết quả thực tế} & Máy chủ từ chối yêu cầu thứ hai, trả về 0 tài nguyên thu hoạch, số dư vàng (Item ID: 00001) của người chơi không bị cộng sai lệch. \\ \hline
\textbf{Trạng thái} & \textbf{PASS} \\ \hline
\end{tabular}
\caption{Kịch bản kiểm thử TC-COLL-02: Ngăn chặn thu hoạch tài nguyên liên tục}
\label{tab:tc_collect_02}
\end{table}

\paragraph{Kịch bản 3: Kiểm thử chức năng Sắp xếp đội hình và Nâng cấp công trình}
Kịch bản kiểm thử kiểm tra các thao tác quản trị đội hình chiến đấu và nâng cấp hạ tầng tông môn:

\begin{table}[H]
\centering
\begin{tabular}{|l|p{10.5cm}|}
\hline
\textbf{Thuộc tính} & \textbf{Nội dung mô tả} \\ \hline
\textbf{Mã kịch bản} & TC-FORM-01 \\ \hline
\textbf{Tên kịch bản} & Sắp xếp đội hình hợp lệ và lưu thành công \\ \hline
\textbf{Mục tiêu} & Đảm bảo người chơi có thể thay đổi vị trí đệ tử trong đội hình $3 \times 3$ và lưu lại trạng thái mới trên cơ sở dữ liệu. \\ \hline
\textbf{Điều kiện} & Người chơi đang sở hữu ít nhất 1 đệ tử trong tông môn. \\ \hline
\textbf{Các bước thực hiện} & 
1. Truy cập giao diện Đội hình. \\
& 2. Chọn đệ tử đang sở hữu và kéo thả vào vị trí số 0 (hàng trước). \\
& 3. Nhấn nút "Lưu đội hình". \\ \hline
\textbf{Dữ liệu đầu vào} & Hero ID: \texttt{501}, Vị trí sắp xếp: \texttt{0} (Hàng trước, ô trái). \\ \hline
\textbf{Kết quả mong đợi} & Client gửi yêu cầu \texttt{SaveFormationRequest} lên World Server, server xác thực đệ tử thuộc sở hữu, cập nhật vị trí vào bảng \texttt{player\_formations} trong PostgreSQL và trả về thành công. \\ \hline
\textbf{Kết quả thực tế} & Cấu hình đội hình được lưu thành công, cơ sở dữ liệu cập nhật chính xác vị trí đệ tử, client hiển thị thông báo "Lưu đội hình thành công". \\ \hline
\textbf{Trạng thái} & \textbf{PASS} \\ \hline
\end{tabular}
\caption{Kịch bản kiểm thử TC-FORM-01: Sắp xếp đội hình hợp lệ}
\label{tab:tc_formation_01}
\end{table}

\begin{table}[H]
\centering
\begin{tabular}{|l|p{10.5cm}|}
\hline
\textbf{Thuộc tính} & \textbf{Nội dung mô tả} \\ \hline
\textbf{Mã kịch bản} & TC-UPGR-01 \\ \hline
\textbf{Tên kịch bản} & Nâng cấp công trình thành công khi đủ tài nguyên \\ \hline
\textbf{Mục tiêu} & Xác minh logic thăng cấp công trình hoạt động chính xác, trừ đúng lượng tài nguyên yêu cầu và cập nhật cấp độ mới. \\ \hline
\textbf{Điều kiện} & Người chơi có đủ lượng gỗ và vàng (vàng (Item ID: 00001)) đáp ứng yêu cầu nâng cấp của công trình. \\ \hline
\textbf{Các bước thực hiện} & 
1. Người chơi chọn công trình Luyện đan phòng (cấp 1). \\
& 2. Nhấn nút "Nâng cấp" trên bảng thông tin chi tiết. \\
& 3. Client gửi yêu cầu \texttt{UpgradeBuildingRequest} lên World Server qua gRPC. \\ \hline
\textbf{Dữ liệu đầu vào} & ID công trình: \texttt{2048}, Cấp hiện tại: 1. Tài sản người chơi: 10.000 vàng (Item ID: 00001) (Yêu cầu: 2.000 vàng (Item ID: 00001)). \\ \hline
\textbf{Kết quả mong đợi} & World Server xác nhận tài sản của người chơi đủ điều kiện, thực hiện trừ 2.000 vàng (Item ID: 00001), nâng cấp công trình lên cấp 2 trong bảng \texttt{player\_buildings} và phản hồi thành công. \\ \hline
\textbf{Kết quả thực tế} & Nâng cấp thành công, cấp công trình tăng lên 2, số dư vàng (Item ID: 00001) của người chơi bị trừ chính xác 2.000 đơn vị, giao diện client hiển thị cấp độ mới của Luyện đan phòng. \\ \hline
\textbf{Trạng thái} & \textbf{PASS} \\ \hline
\end{tabular}
\caption{Kịch bản kiểm thử TC-UPGR-01: Nâng cấp công trình thành công}
\label{tab:tc_upgrade_01}
\end{table}

\subsubsection{Tổng kết kết quả kiểm thử}
Hệ thống đã thực hiện tổng cộng 25 trường hợp kiểm thử chức năng tự động (Unit Test) và thủ công (Manual Test) bao quát tất cả các tính năng cơ bản. Kết quả ghi nhận đạt tỷ lệ thành công 100\% (25/25 trường hợp kiểm thử đạt trạng thái PASS), không phát hiện lỗi logic nghiêm trọng trong các tiến trình trừ/cộng tài nguyên và thăng cấp nhân vật.

\section{Triển khai}
\label{sec:trien_khai}

\subsection{Mô hình triển khai hệ thống}
Để đáp ứng yêu cầu vận hành liên tục 24/7 với độ sẵn sàng dịch vụ tối thiểu đạt 99.5\%, hệ thống máy chủ được ảo hóa hoàn toàn bằng Docker và điều phối thông qua Docker Compose. Toàn bộ hạ tầng được thiết lập chạy trên hệ thống máy chủ đám mây ảo (Cloud VPS) cấu hình tiêu chuẩn: CPU 4 nhân (vCPU), 8 GB RAM, ổ cứng SSD 80 GB chạy hệ điều hành Ubuntu Server 22.04 LTS.

\begin{figure}[H]
\centering
\begin{tikzpicture}[node distance=1.1cm and 1.2cm, every node/.style={font=\small, align=center}]
  % Internet Client
  \node[draw, rectangle, fill=orange!10, minimum width=2.5cm, minimum height=0.8cm] (client) {Client (Unity 6)};
  
  % Nginx Gateway
  \node[draw, rectangle, fill=teal!20, minimum width=3.5cm, minimum height=0.8cm, below=of client] (nginx) {Nginx Reverse Proxy \\ (Port 80/443, SSL, Routing)};
  
  % Internal Gateway Server
  \node[draw, rectangle, fill=teal!10, minimum width=3.5cm, minimum height=0.8cm, below=of nginx] (gateway) {Gateway Server (Go) \\ (Port 50051)};

  % Microservices inside Docker Network
  \node[draw, rectangle, fill=green!10, minimum width=2.2cm, minimum height=0.8cm, below=of gateway] (world) {World Server \\ (Port 50053)};
  \node[draw, rectangle, fill=green!10, minimum width=2.2cm, minimum height=0.8cm, left=of world] (login) {Login Server \\ (Port 50052)};
  \node[draw, rectangle, fill=green!10, minimum width=2.2cm, minimum height=0.8cm, right=of world] (combat) {Combat Server \\ (Port 50054)};

  % Database Connections
  \node[draw, cylinder, shape border rotate=90, fill=purple!10, minimum width=1.5cm, minimum height=1.4cm, below=1.2cm of world, aspect=0.15] (pg_game) {PostgreSQL \\ (Game Shard)};
  \node[draw, cylinder, shape border rotate=90, fill=purple!10, minimum width=1.5cm, minimum height=1.4cm, below=1.2cm of login, aspect=0.15] (pg_global) {PostgreSQL \\ (Global)};
  \node[draw, cylinder, shape border rotate=90, fill=purple!10, minimum width=1.5cm, minimum height=1.4cm, below=1.2cm of combat, aspect=0.15] (redis) {Redis};

  % Routing lines
  \draw[->, >=stealth, thick] (client) -- (nginx) node[midway, right] {gRPC over TLS};
  \draw[->, >=stealth, thick] (nginx) -- (gateway) node[midway, right] {Forward HTTP/2};
  
  \draw[->, >=stealth, thick] (gateway) -- (login);
  \draw[->, >=stealth, thick] (gateway) -- (world);
  \draw[->, >=stealth, thick] (gateway) -- (combat);

  \draw[->, >=stealth, dashed] (login) -- (pg_global);
  \draw[->, >=stealth, dashed] (world) -- (pg_game);
  \draw[->, >=stealth, dashed] (combat) -- (redis);
\end{tikzpicture}
\caption{Mô hình triển khai máy chủ vi dịch vụ sử dụng Docker và Nginx}
\label{fig:deployment_model}
\end{figure}

\subsection{Cấu hình mạng và Cân bằng tải (Reverse Proxy)}
Hệ thống sử dụng \textbf{Nginx Reverse Proxy} làm cổng tiếp nhận ngoài cùng để thực hiện:
\begin{itemize}
    \item \textbf{Đóng gói SSL (SSL Termination)}: Giải mã kết nối HTTPS/gRPC sử dụng chứng chỉ mã hóa trước khi chuyển gói tin dưới dạng HTTP/2 thuần túy vào mạng nội bộ Docker.
    \item \textbf{Định tuyến dựa trên dịch vụ (Path-based Routing)}: Nginx phân tích tên dịch vụ gRPC trong đường dẫn yêu cầu để định tuyến chính xác. Ví dụ, các gói tin yêu cầu gọi dịch vụ \texttt{/pb.LoginService/} sẽ được tự động chuyển tiếp tới container của Login Server, trong khi các cuộc gọi \texttt{/pb.WorldService/} sẽ chuyển tiếp tới World Server tương ứng.
\end{itemize}

\subsection{Cơ chế Roll Server điều phối phân vùng}
Để xử lý vấn đề quá tải số lượng người chơi kết nối đồng thời vào một máy chủ đơn lẻ, hệ thống áp dụng cơ chế tự động điều phối phân vùng (Roll Server gateway):
\begin{enumerate}
    \item Khi khởi động trò chơi, Client gửi yêu cầu truy vấn danh sách máy chủ khả dụng (\texttt{GetZoneList}) tới Gateway Server.
    \item Gateway Server tiến hành truy cập bảng \texttt{zones} trong PostgreSQL Global và thực hiện lấy số liệu hiệu năng (RAM, CPU, số lượng kết nối đang hoạt động) của các Zone đang trực tuyến.
    \item Gateway tự động phân loại trạng thái tải của từng Zone thành các mức: \textit{Bình thường (normal)}, \textit{Đông đúc (crowded)}, hoặc \textit{Bảo trì (maintenance)}.
    \item Danh sách máy chủ kèm địa chỉ URL định tuyến cụ thể (ví dụ \url{127.0.0.1:50053}) được trả về Client. Nếu người chơi chọn chế độ tự động điều phối hoặc chọn một máy chủ đã đầy, Gateway sẽ tự động định hướng yêu cầu kết nối của Client sang phân vùng máy chủ còn trống (Roll Server) để cân bằng tải trọng kết nối, ngăn ngừa nguy cơ sập cục bộ một phân vùng do quá tải.
\end{enumerate}

\subsection{Đánh giá kết quả triển khai thử nghiệm}
Để đánh giá năng lực chịu tải thực tế, đồ án đã triển khai chương trình giả lập người chơi ảo (\textit{Bot Simulator}) viết bằng ngôn ngữ Go. Chương trình này chạy song song 5.000 Goroutines tương ứng với 5.000 người chơi ảo (CCU) kết nối trực tiếp đến hệ thống gRPC Gateway.

\subsubsection{Cấu hình giới hạn tần suất yêu cầu (Rate Limiting)}
Để ngăn chặn các cuộc tấn công từ chối dịch vụ (DoS) và kiểm soát lưu lượng kết nối của từng người chơi, Gateway Server được trang bị bộ lọc trung gian (\textit{UnaryInterceptor}) giới hạn tần suất yêu cầu (\textit{Rate Limiter}) dựa trên thuật toán Token Bucket. Hệ thống trích xuất địa chỉ IP của Client bằng cách tách thông tin cổng qua hàm \texttt{net.SplitHostPort} để quản lý chính xác từng nguồn kết nối độc lập.
\begin{itemize}
    \item \textbf{Cấu hình mặc định:} Tần suất trung bình (\textit{RPS}) là 100 req/s, dung lượng chứa tối đa (\textit{Burst}) là 200.
    \item \textbf{Khả năng tùy biến:} Các thông số trên được thiết lập linh hoạt qua biến môi trường (\texttt{RATE\_\allowbreak LIMIT\_\allowbreak ENABLED}, \texttt{RATE\_\allowbreak LIMIT\_\allowbreak RPS}, \texttt{RATE\_\allowbreak LIMIT\_\allowbreak BURST}), cho phép quản trị viên tắt hoặc tăng giới hạn này khi cần chạy các bài kiểm thử tải quy mô lớn từ một địa chỉ IP duy nhất.
\end{itemize}

\subsubsection{Kết quả thực nghiệm}
Trong bài kiểm thử tải trọng, quá trình đo đạc được tiến hành với các mức số lượng người chơi giả lập (CCU) khác nhau bao gồm 100 bots, 1.000 bots và 5.000 bots. Cơ chế giới hạn tần suất được tạm thời vô hiệu hóa (\texttt{RATE\_\allowbreak LIMIT\_\allowbreak ENABLED=false}) trong suốt thời gian chạy để tránh việc chặn lưu lượng giả lập từ một địa chỉ IP duy nhất. Kết quả ghi nhận chi tiết như sau:

\begin{table}[H]
\centering
\caption{So sánh thông số hiệu năng hệ thống dưới các mức tải CCU khác nhau (Thời gian test: 5 phút)}
\label{tab:load_test_comparison}
\resizebox{\textwidth}{!}{
\begin{tabular}{|l|c|c|c|}
\hline
\textbf{Thông số đo lường} & \textbf{Tải nhẹ (100 CCU)} & \textbf{Tải trung bình (1.000 CCU)} & \textbf{Tải cao (5.000 CCU)} \\ \hline
Tổng số yêu cầu đã gửi & 11.485 & 103.655 & 110.731 \\ \hline
Tốc độ xử lý (TPS) & 38,28 req/s & 345,51 req/s & 362,19 req/s \\ \hline
Tỷ lệ yêu cầu thành công & 100,0\% & 100,0\% & 99,80\% (110.507 req) \\ \hline
Độ trễ phản hồi trung bình & 2 ms & 4 ms & 5 ms \\ \hline
Độ trễ thấp nhất (Min) & 1 ms & 1 ms & 1 ms \\ \hline
Độ trễ cao nhất (Max) & 45 ms & 98 ms & 313 ms \\ \hline
Sử dụng CPU máy chủ (Go) & 2,5\% & 12,0\% & 64,96\% \\ \hline
Sử dụng RAM máy chủ & 20 MB & 35 MB & 98 MB \\ \hline
\end{tabular}
}
\end{table}

\textbf{Nhận xét và phân tích:}
\begin{itemize}
    \item Ở mức tải nhẹ và tải trung bình (100 CCU và 1.000 CCU), hệ thống hoạt động hoàn hảo với tỷ lệ yêu cầu thành công tuyệt đối (100\%).
    \item Khi tải tăng mạnh lên mức tải đỉnh (5.000 CCU) và duy trì liên tục trong 5 phút, tốc độ xử lý (TPS) đạt mức ổn định ở 362,19 req/s. Sự sụt giảm nhỏ của tỷ lệ yêu cầu thành công xuống còn 99,80\% (tương ứng 224 request bị thất bại) xảy ra chủ yếu trong giai đoạn đầu đẩy tải, xuất phát từ lỗi cạn kiệt cổng kết nối cục bộ ở hệ điều hành Windows (\textit{TCP Port Exhaustion/Socket Refusal}), khiến một số bot bị chặn từ ngoài cửa. Tuy nhiên, hệ thống máy chủ Go vẫn giữ được độ trễ trung bình thấp (chỉ 5 ms) và tài nguyên ổn định xuyên suốt quá trình kiểm thử, chứng minh cấu trúc xử lý đồng thời không hề bị nghẽn cổ chai (bottleneck) tại lớp ứng dụng (Server Layer).
\end{itemize}

\vspace{0.5cm}
Chương này đã làm rõ toàn bộ quy trình thiết kế kiến trúc hệ thống tổng thể, thiết kế cơ sở dữ liệu chi tiết, các kịch bản kiểm thử chức năng nghiêm ngặt và giải pháp triển khai thực tế. Những kết quả kiểm nghiệm hiệu năng đạt được là cơ sở vững chắc để đồ án đề xuất các giải pháp tối ưu hóa chuyên sâu sẽ được trình bày chi tiết trong Chương 5.


\end{document}
